\section{System and Scenario}
\label{sec:system}

All results in this paper are for the simulated scenario described here; it is the same scenario in every experiment except where a section says otherwise.
% src: A3; configs/demo.yaml; D20-D22, D29-D32, D34, D41; F10

\paragraph*{Network and mobility}
Mobility is synthetic, on a Manhattan grid of arterials ($1600$\,m square, three blocks per side), not generated by a traffic simulator.
% src: A3; configs/demo.yaml (road.extent_m 1600, blocks_x 3, blocks_y 3); L7
The network has 20 RSUs and 60 vehicles, and each RSU is homed to one of 4 backhaul segments.
% src: A3; configs/demo.yaml (num_rsus 20, num_vehicles 60, num_backhaul_segments 4)
The network is rebuilt as a graph at every one-second step: nodes are RSUs and vehicles, edges are vehicle--RSU coverage links and RSU--RSU coordination links, and a single edge type carries a \texttt{same\_segment} flag that is $1$ when both RSU endpoints share a backhaul segment.
% src: A3 (same_segment, L6); PROJECT_SPEC 5.3; configs/demo.yaml (dt_s 1.0)
Each run has 1200 steps.
% src: configs/demo.yaml (num_steps 1200)

\paragraph*{Tasks and execution}
Each vehicle generates tasks at a rate of $0.1$\,Hz, of three types (light, medium, heavy) with shares $0.5$, $0.35$ and $0.15$; the deadline is a communication allowance plus a slack over the compute time on an idle reference node, capped at $1$\,s.
% src: A3 (three task types, closed-form execution model); configs/demo.yaml (tasks.rate_per_vehicle_hz 0.1, shares [0.5,0.35,0.15], deadline_cap_s 1.0)
Execution is closed-form: an RSU's runtime depends on its true load and a mean-one lognormal scatter.
% src: A3 (closed-form execution model); configs/demo.yaml (execution.noise_sigma 0.25)

\paragraph*{Degradation}
Degradation is injected per backhaul segment, not per node, so that correlated failure is a structural property of the topology.
% src: A3 (L5)
A degraded RSU loses compute capacity and silently drops tasks; onset is shared by all degraded RSUs at step $400$ and severity ramps linearly over $300$ steps, to $0.6$ of capacity lost and a drop probability of $0.15$ at full degradation.
% src: A3 (shared onset at step 400, linear ramp over 300 steps); configs/demo.yaml (severity 0.6, drop_prob 0.15)
The number of degraded RSUs is $K=\lfloor 0.5+f\cdot 20\rfloor$ for degraded fraction $f$, and the failure correlation $\rho\in[0,1]$ controls how concentrated they are on one segment: $\rho=0$ picks degraded RSUs uniformly, and $\rho=1$ fills one segment before opening the next.
% src: A3 (rho interpolates concentration, D29); DECISIONS D29
The realised correlation is recorded and reported beside nominal $\rho$.
At a degraded fraction of 5\% only one RSU is degraded, so there is no correlated failure at that point; at 10\% every draw is binary.
% src: D43 section 1 and D56 section 1 (K = 1 at 5%, K = 2 at 10%); F10

\paragraph*{What the trust signal is}
Nodes advertise load and latency, and after a task completes or times out, the observed outcome is compared with what was advertised.
The \emph{discrepancy}, not the advertised value, feeds two behavioural features, \texttt{success\_ewma} and \texttt{latency\_dev}.
% src: A3 (L8); PROJECT_SPEC 5.1
An RSU row also carries \texttt{load}, \texttt{queue\_depth}, \texttt{cert\_valid} and \texttt{uptime\_stability}.
In this scenario \texttt{cert\_valid} and \texttt{uptime\_stability} are constant for every RSU at every step, so the informative inputs of an RSU are its advertised load and queue depth and the two discrepancy features.
% src: F8, D40 (placeholders: 1.0 for every RSU at every step); A4 (cert_valid constant). NOT an A-entry: see S12 report.
The injected behaviour class (reliable, degraded, compromised) lives in a separate sealed package that the training path cannot import; a test enforces this, and it is read only by evaluation code.
% src: A2 (sealed, tested); L4, D35

\paragraph*{Seeds, draws and the unit of analysis}
A seed fixes the topology (including backhaul segment assignment), traffic and execution noise.
An \emph{injection draw} redraws which RSUs degrade for that seed's fixed traffic.
% src: D39, D44 (training and evaluation share a seed's traffic and differ in injection draws; A19)
There are 5 seeds; the unit of analysis is the seed, and every interval in this paper is a 95\% Student-$t$ interval over the 5 per-seed values (4 degrees of freedom).
% src: A16 (unit of analysis, intervals); D43
Mean $\pm$ standard deviation is across seeds.
